% Part:sets-functions-relations
% Chapter: size-of-sets
% Section: equinumerous-sets

\documentclass[../../../include/open-logic-section]{subfiles}

\begin{document}

\olfileid{sfr}{siz}{equ}

\olsection{Gelijkmachtigheid}

We hebben een intuïtief begrip van de ``grootte'' van verzamelingen, dat voor
eindige verzamelingen goed werkt. Maar hoe zit het met oneindige verzamelingen?
Voor een formele vergelijking van de grootte van twee verzamelingen van
\emph{willekeurige} grootte kunnen we het best eerst vastleggen wanneer
verzamelingen even groot zijn. Frege schrijft:
\begin{quote}
  Als een ober er zeker van wil zijn dat hij precies evenveel messen als
  borden op tafel heeft gelegd, hoeft hij geen van beide te tellen. Hij kan
  eenvoudig rechts van elk bord een mes leggen, waarbij elk mes op tafel
  rechts van een bord ligt. De borden en messen zijn zo één-op-één aan elkaar
  gekoppeld, en wel door diezelfde ruimtelijke betrekking. \citep[\S70]{Frege1884}
\end{quote}
Het inzicht uit deze passage kunnen we in een formele definitie vastleggen:

\begin{defn}\ollabel{comparisondef}
  $A$ is \emph{gelijkmachtig} met $B$, geschreven als $\cardeq{A}{B}$,
  precies dan wanneer er !!a{bijection} $f \colon A \to B$ bestaat.
\end{defn}

\begin{prop}\ollabel{equinumerosityisequi}
Gelijkmachtigheid is een equivalentierelatie.
\end{prop}

\begin{proof} 
We moeten aantonen dat gelijkmachtigheid reflexief, symmetrisch en transitief
is. Laat $A, B$ en $C$ verzamelingen zijn.

\emph{Reflexiviteit.} De identiteitsfunctie $\Id{A} \colon A \to A$, met
$\Id{A} (x) = x$ voor alle $x \in A$, is !!a{bijection}. Dus geldt
$\cardeq{A}{A}$.

\emph{Symmetrie.} Stel dat $\cardeq{A}{B}$, dus dat er !!a{bijection}
$f\colon A \to B$ bestaat. Omdat $f$ !!{bijective} is, bestaat zijn inverse
$f^{-1}$ en is deze eveneens !!{bijective}. Bijgevolg is
$f^{-1}\colon B \to A$ !!a{bijection}, zodat $\cardeq{B}{A}$.

\emph{Transitiviteit.} Stel dat $\cardeq{A}{B}$ en $\cardeq{B}{C}$, dus dat
er !!{bijection}s $f\colon A \to B$ en $g\colon B \to C$ bestaan. Dan is de
samenstelling $\comp{f}{g}\colon A \to C$ !!{bijective}, zodat
$\cardeq{A}{C}$.
\end{proof}

\begin{prop}
Als $\cardeq{A}{B}$, dan is $A$ precies dan !!{enumerable} als $B$ dat is.
\end{prop}

\begin{editorial}
Het volgende bewijs gebruikt \olref[enm]{defn:enumerable} als
\olref[enm]{sec} is opgenomen en anders \olref[enm-alt]{defn:enumerable}.
\end{editorial}

\begin{proof}
Stel dat $\cardeq{A}{B}$. Er bestaat dan !!a{bijection} $f \colon A
\to B$. Stel bovendien dat $A$ !!{enumerable} is.
\oliflabeldef{sfr:siz:enm:defn:enumerable}{
  Dan geldt $A = \emptyset$ of er bestaat een surjectieve functie
  $g\colon \PosInt \to A$. Als $A = \emptyset$, dan is ook $B = \emptyset$
  (anders zou er !!a{element}~$y \in B$ zijn, maar geen $x \in
  A$ waarvoor $f(x) = y$). Als daarentegen $g\colon \PosInt \to A$
  !!{surjective} is, dan is $\comp{g}{f} \colon \PosInt \to B$
  !!{surjective}. Neem om dit te zien een willekeurige $y \in B$. Omdat $f$
  !!{surjective} is, bestaat er een $x \in A$ waarvoor $f(x) = y$. Omdat
  $g$ !!{surjective} is, bestaat er een $n \in \PosInt$ waarvoor $g(n) =
  x$. Bijgevolg
\[
(\comp{g}{f})(n) = f(g(n)) = f(x) = y
\]
en dus is $\comp{g}{f}$ !!{surjective}. Daarmee is $\comp{g}{f}$ een
opsomming van~$B$, zodat $B$~!!{enumerable} is.}
{
  Dan geldt $A = \emptyset$ of er bestaat !!a{bijection}~$g$ met bereik
  $A$ en met als domein $\Nat$ of een beginstuk van de natuurlijke getallen. Als
  $A = \emptyset$, dan is ook $B = \emptyset$ (anders zou er een~$y \in B$
  zijn zonder $x \in A$ waarvoor $f(x) = y$). Neem dus aan dat we zo'n
  !!{bijection}~$g$ hebben. Dan is $\comp{g}{f}$ !!a{bijection} met bereik
  $B$ en hetzelfde domein als~$g$ (dus $\Nat$ of een beginstuk daarvan),
  zodat $B$ !!{enumerable} is.}

Als $B$ !!{enumerable} is, volgt op dezelfde manier dat $A$ !!{enumerable}
is, nu met de !!{bijection} $f^{-1}\colon B \to A$ in plaats van~$f$.
\end{proof}

\begin{prob}
Toon aan dat uit $\cardeq{A}{C}$, $\cardeq{B}{D}$ en $A \cap B =
C \cap D = \emptyset$ volgt dat $\cardeq{A \cup B}{C \cup D}$.
\end{prob}

\begin{prob}
Toon aan dat als $A$ oneindig en !!{enumerable} is, dan
$\cardeq{A}{\Nat}$.
\end{prob}
\end{document}
